% Part: first-order-logic
% Chapter: tableaux
% Section: provability-quantifiers

\documentclass[../../../include/open-logic-section]{subfiles}

\begin{document}

\olfileid{fol}{tab}{qpr}
\olsection{\usetoken{S}{derivability} تے مسوّر (Derivability and the Quantifiers)}

\begin{explain}
  تمامیت دے قضیے دا تقاضا ایہ وی اے کہ تابلو دے قاعدے~$\Proves$
  بارے اوہ حقائق دیندے ہون جیہڑے ایس حصے وچ قائم کیتے گئے نیں۔
\end{explain}

\begin{thm}
\ollabel{thm:strong-generalization} جے $c$ اک ایہو جہا مستقل (constant) اے جیہڑا
$\Gamma$ یا~$!A(x)$ وچ نہیں آؤندا تے $\Gamma \Proves !A(c)$، تاں $\Gamma
\Proves \lforall[x][!A(x)]$۔
\end{thm}

\begin{proof}
فرض کرو $\Gamma \Proves !A(c)$، یعنی $!B_1$، \dots، $!B_n
\in \Gamma$ موجود نیں تے ہِٹھلے سیٹ لئی اک بند !!{tableau} اے:
\begin{align*}
\{ \sFmla{\False}{!A(c)}, &
\sFmla{\True}{!B_1}, \dots, \sFmla{\True}{!B_n} \}.
\intertext{سانوں وکھانا اے کہ ہِٹھلے سیٹ لئی وی اک بند !!{tableau} موجود اے:}
\{ \sFmla{\False}{\lforall[x][!A(x)]}, &
\sFmla{\True}{!B_1}, \dots, \sFmla{\True}{!B_n} \}.
\end{align*}
ایہ بند !!{tableau} لو تے پہلا مفروضہ
$\sFmla{\False}{\lforall[x][!A(x)]}$ نال بدلو؛ فیر سارے
مفروضیاں توں بعد $\sFmla{\False}{!A(c)}$ شامل کرو۔
\begin{center}
\begin{tableau}{not line numbering}
  [\sFmla{\False}{\formula{A}(c)}
    [\sFmla{\True}{\formula{B}_1}
      [\vdots
      [\sFmla{\True}{\formula{B}_n}, 
        [][][]]]]]
\end{tableau}\qquad
\begin{tableau}{not line numbering}
  [\sFmla{\False}{\lforall[x][\formula{A}(x)]}
    [\sFmla{\True}{\formula{B}_1}
      [\vdots
        [\sFmla{\True}{\formula{B}_n},
          [\sFmla{\False}{\formula{A}(c)}
        [][][]]]]]]
\end{tableau}
\end{center}
ایہ تابلو ہن وی بند اے، کیوں جے جیہڑے سارے
!!{sentence}s پہلے مفروضیاں دے طور تے دستیاب سن، اوہ ہن وی
!!{tableau} دے شروع وچ موجود نیں۔ شامل کیتی لکیر
$\TRule{\False}{\lforall}$ دی درست ورتوں نال ملی اے، کیوں جے
!!{constant}~$c$، $!B_1$، \dots، $!B_n$ یا
$\lforall[x][!A(x)]$ وچ نہیں آؤندا؛ یعنی نویں !!{tableau}
وچ شامل کیتی لکیر توں اُتّے کہیں وی نہیں آؤندا۔
\end{proof}

\begin{prop}
\ollabel{prop:provability-quantifiers}
\begin{tagenumerate}{prvEx,prvAll}
\tagitem{prvEx}{$!A(t) \Proves \lexists[x][!A(x)]$.}{}
\tagitem{prvAll}{$\lforall[x][!A(x)] \Proves !A(t)$.}{}
\end{tagenumerate}
\end{prop}

\begin{proof}
\begin{tagenumerate}{prvEx,prvAll}
\tagitem{prvEx}{$\sFmla{\False}{\lexists[x][!A(x)]}, \sFmla{\True}{!A(t)}$
  لئی اک بند !!{tableau} ایہ اے:
  \begin{oltableau}
    [\sFmla{\False}{\lexists[x][\formula{A}(x)]}, just = \TAss
      [\sFmla{\True}{\formula{A}(t)}, just = \TAss
        [\sFmla{\False}{\formula{A}(t)}, just = {\TRule{\False}{\lexists}[1]}, close]]]
    \end{oltableau}}{}
% PNBTAB-012: source line 79's trailing comma is punctuation, not part of the signed premise.
\tagitem{prvAll}{$\sFmla{\False}{\formula{A}(t)}, \sFmla{\True}{\lforall[x][!A(x)]}$
  لئی اک بند !!{tableau} ایہ اے:
  \begin{oltableau}
    [\sFmla{\False}{\formula{A}(t)}, just = \TAss
      [\sFmla{\True}{\lforall[x][\formula{A}(x)]}, just = \TAss
        [\sFmla{\True}{\formula{A}(t)}, just = {\TRule{\True}{\lforall}[2]}, close]]]
    \end{oltableau}}{}
\end{tagenumerate}
\end{proof}

\end{document}
